\section*{Chapitre \ref{ch:rs:fundamental-analyst-and-event-features} --- Prédicteurs fondamentaux, d'analystes et d'événements}

\begin{solution}{exo:rs:fundamental-analyst-and-event-features:1}
$\$4.80 - \$3.30 = \$1.50$. Elle n'est connue pour la première fois qu'au dépôt du 10-K, 55 jours après le 30 juin~: le 24 août (ou à la publication des résultats, si
le magasin de données est indexé sur les publications).
\end{solution}

\begin{solution}{exo:rs:fundamental-analyst-and-event-features:2}
Quarante jours après le 30 septembre~: le 9 novembre. Soixante jours après le 31 décembre, dans une année de 365 jours~: le premier mars. (Une échéance qui tombe un week-end
ou un jour férié est reportée au jour ouvré suivant.)
\end{solution}

\begin{solution}{exo:rs:fundamental-analyst-and-event-features:3}
$1.42/0.355 = 4$~: une division d'actions de quatre pour un a rééchelonné l'historique. Le magasin conserve 1.42~\$ avec la date de son premier dépôt comme \omterm{def:rs:point-in-time-data-and-the-biases:bitemporal}{temps de connaissance}, et 0.355~\$
avec celle du dépôt ultérieur, et enregistre la division comme un événement doté de sa propre date~; une requête ramène les valeurs à la base par action de son instant de décision.
\end{solution}

\begin{solution}{exo:rs:fundamental-analyst-and-event-features:4}
Les différences saisonnières valent 0.04, 0.02, 0.06, 0.02, 0.02, 0.04, $-0.02$ et 0.18. Les six précédant la dernière (0.02, 0.06, 0.02, 0.02, 0.04,
$-0.02$) ont un écart-type d'échantillon de 0.0266, donc le SUE du dernier trimestre vaut $0.18/0.0266 = 6.8$.
\end{solution}

\begin{solution}{exo:rs:fundamental-analyst-and-event-features:5}
La fenêtre allant du jour 1 au jour 21 après la fin du trimestre contient l'annonce lorsque celle-ci tombe sur l'un de ces 21 jours parmi les 63 possibles~: un
tiers (0.330 mesuré). Dans ces fenêtres, le prédicteur voit le saut de l'annonce, dont l'ampleur écrase la dérive que mesure la bonne clé~; il
suffit qu'un tiers des fenêtres porte un saut plusieurs fois supérieur à la dérive d'un trimestre pour multiplier l'IC par huit.
\end{solution}

\begin{solution}{exo:rs:fundamental-analyst-and-event-features:6}
Un \omterm{def:rs:point-in-time-data-and-the-biases:lookahead}{biais d'anticipation} gagne ce que le cours fait entre la fin de la période et le \omterm{def:rs:point-in-time-data-and-the-biases:bitemporal}{temps de connaissance}, à cause du chiffre caché. Une \omterm{def:rs:fundamental-analyst-and-event-features:sue}{surprise de résultats} fait bouger le
cours lors de sa publication, dans l'intervalle~; le dépôt d'une valeur comptable le fait peu bouger, et le book-to-price évolue de toute façon lentement, si bien que le rendement de la fenêtre est
presque indépendant du choix entre les deux clés. Le book-to-price gonflerait si la valeur comptable contenait une information à laquelle le cours a réagi lors du dépôt
(une forte dépréciation), ou si le cours de fin de période utilisé dans le ratio était remplacé par un cours postérieur (un \omterm{def:rs:point-in-time-data-and-the-biases:lookahead}{biais d'anticipation} sur les cours).
\end{solution}

\begin{solution}{exo:rs:fundamental-analyst-and-event-features:7}
\texttt{rs\_fundamentals.release\_to\_filing}~: Microsoft a déposé chacun de ses 10-Q le jour de la publication, Procter \& Gamble 95~\% d'entre eux~; Coca-Cola 7~\% et
Intel 10~\%. Pour les autres, les délais médians entre la publication et le dépôt sont d'un jour (Apple, Intel), de deux jours (Coca-Cola, Procter \& Gamble), de cinq jours (Exxon Mobil),
de sept jours (Home Depot), de 14 jours (Johnson \& Johnson, Nike) et de 16 jours (Walmart).
\end{solution}

\begin{solution}{exo:rs:fundamental-analyst-and-event-features:8}
C'est l'historique le plus exact et le mauvais pour un backtest~: chaque valeur est celle d'aujourd'hui, si bien qu'une décision passée voit les \omterm{def:rs:point-in-time-data-and-the-biases:vintage}{retraitements} (le
BPA de Microsoft pour l'exercice 2017, 3.25~\$ au lieu des 2.71~\$ connus jusqu'en août 2018) et les rééchelonnements dus aux divisions d'actions qu'elle ne pouvait pas connaître. La recherche a besoin des valeurs telles
qu'elles ont été publiées pour la première fois, avec chaque version ultérieure et sa date.
\end{solution}

\begin{solution}{pb:rs:fundamental-analyst-and-event-features:1}
\begin{enumerate}
\item 542 fins de période~; 541 sont suivies d'une publication dans les 90 jours.
\item Trimestres~: 23 jours jusqu'à la publication, 29 jusqu'au 10-Q. Exercices annuels~: 26 jours jusqu'à la publication, 48 jusqu'au 10-K.
\item 41 et 60 jours~; les échéances sont de 40 jours (reportées au jour ouvré suivant lorsqu'elles tombent un jour férié ou un week-end) et de 60 jours pour les grands déposants.
\item 16~\% le jour même~; l'écart médian est de 7 jours.
\item 31, dont 24 déterminées par un ratio exact.
\item Division par sept le 6 juin 2014 et par quatre le 28 août 2020~; les valeurs par action antérieures à chacune ont été ensuite publiées divisées par 7 et par 4 dans
les dépôts qui les reprenaient.
\item Publié d'abord à 2.71~\$, retraité à 3.25~\$ dans les dépôts de l'exercice 2018~: Microsoft a adopté une nouvelle norme sur le chiffre d'affaires en appliquant la
méthode rétrospective intégrale, qui a retraité chaque période antérieure présentée.
\item 28.82~\$ à partir des trimestres publiés d'abord (deux avant la division, deux après), contre un BPA annuel de 6.45~\$.
\item Tels que publiés d'abord avec leurs \omterm{def:rs:point-in-time-data-and-the-biases:bitemporal}{temps de connaissance}, les divisions étant stockées comme événements et ramenées lors de la requête à la
base par action de la date de décision.
\item 31 jours de bourse, soit environ six semaines.
\item 0.020 (erreur type 0.006) indexé sur l'annonce~; 0.167 (0.006) indexé sur la fin de période.
\item Un tiers (0.330).
\item 0.029 (0.038) indexé sur la fin de période~; $-0.005$ (0.023) indexé sur le dépôt, avec le cours de fin de période.
\item L'IC de chaque trimestre est dominé par le rendement du facteur valeur dans sa fenêtre, et rien ne se produit le jour du dépôt qui fasse dépendre les
rendements des deux fenêtres du ratio de façon différente~: le gonflement attendu est faible et le bruit important.
\item \textbf{Résultat nommé.} Dans \texttt{firm.synthmkt}, indexer la \omterm{def:rs:fundamental-analyst-and-event-features:sue}{surprise de résultats} sur la fin de période au lieu de son annonce multiplie
son IC par huit (0.020 à 0.167)~; indexer le book-to-price sur la fin de période au lieu de son dépôt le modifie de moins que son erreur type.
\item Les jours écoulés entre la publication et le dépôt (une semaine à la médiane dans les données EDGAR)~: pour un prédicteur de dérive, le début de la dérive.
\item La clé de fin de période, de loin, pour un prédicteur SUE~: il porte sur l'annonce elle-même. Un historique retraité déforme la taille de certaines surprises
mais pas leur moment.
\item Recalculer l'IC du prédicteur avec l'entrée décalée au \omterm{def:rs:point-in-time-data-and-the-biases:bitemporal}{temps de connaissance}~; tracer l'IC en fonction du décalage, car une fuite se manifeste par une
rupture brutale au véritable \omterm{def:rs:fundamental-analyst-and-event-features:fiscal}{délai de publication}~; vérifier le moment des rendements réalisés autour des annonces.
\item Chaque valeur avec sa période, son \omterm{def:rs:point-in-time-data-and-the-biases:bitemporal}{temps de connaissance} initial, chaque version ultérieure avec son temps, et les opérations sur le capital consignées comme événements datés.
\item L'instant de sa première publication, à la publication des résultats ou, pour une base alimentée par les dépôts, au dépôt, jamais la fin de la période qu'elle décrit.
\end{enumerate}
\end{solution}

\begin{solution}{iq:rs:fundamental-analyst-and-event-features:1}
Des données stockées avec l'instant où chaque valeur est devenue connue, et interrogées à un instant de décision, si bien qu'un backtest ne voit que ce qui était connu à cet instant. Cela importe
parce que les chiffres comptables sont publiés des semaines après leurs périodes et sont ensuite retraités ou rééchelonnés~; une base de valeurs courantes indexée sur
les fins de période fuit de deux façons.
\end{solution}

\begin{solution}{iq:rs:fundamental-analyst-and-event-features:2}
Calculer la différence des chiffres cumulés depuis le début de l'exercice~: le deuxième trimestre vaut les six mois moins le premier, le quatrième l'exercice entier moins les neuf mois~; chaque
trimestre dérivé est connu quand sa composante la plus tardive l'est. Sommer les quatre trimestres consécutifs tels que connus à l'instant de décision, sur une même base par action pour les chiffres par action.
\end{solution}

\begin{solution}{iq:rs:fundamental-analyst-and-event-features:3}
La tendance des rendements, après une publication de résultats, à se poursuivre dans le sens de la surprise (Foster, Olsen et Shevlin, 1984). Prédicteur~: le SUE
contre une marche aléatoire saisonnière ou la prévision consensuelle, indexé sur la publication, classé parmi les actions ayant annoncé récemment, conservé pendant des semaines~; vérifier
la date d'indexation et neutraliser la taille.
\end{solution}

\begin{solution}{iq:rs:fundamental-analyst-and-event-features:4}
La date de dépôt est la bonne~: la fin de période utilise des chiffres avant qu'ils ne soient connus. L'amélioration est grande lorsque la publication du chiffre fait bouger le
cours (résultats) et faible lorsqu'elle ne le fait pas (une valeur comptable qui évolue lentement)~; sur le marché simulé du chapitre, le book-to-price a bougé de moins que son
erreur type, tandis que la surprise a été multipliée par huit. Un tiers est un avertissement dans les deux cas.
\end{solution}

\begin{solution}{iq:rs:fundamental-analyst-and-event-features:5}
Horodatages (date de la prévision comparée à la date où le fournisseur l'a enregistrée), révisions qui écrasent l'historique, \omterm{def:rs:point-in-time-data-and-the-biases:survivorship}{biais de survie} des analystes et des entreprises suivies,
choix de la couverture en fonction de la performance attendue, prévisions périmées dans la prévision consensuelle, et optimisme qui varie avec les incitations de l'analyste.
\end{solution}

\begin{solution}{iq:rs:fundamental-analyst-and-event-features:6}
Un entrepôt bitemporel de faits indexé par (entité, champ, période), avec chaque version et son \omterm{def:rs:point-in-time-data-and-the-biases:bitemporal}{temps de connaissance}, les valeurs conservées telles que publiées sur leur propre
base par action~; une table distincte des événements de division avec leurs dates~; une requête (entité, champ, période, instant de décision) qui prend la dernière version connue à ce moment et
la multiplie par le produit des facteurs de division entre la base de la version et celle de la date de décision, toutes deux prises comme connues à l'instant de décision.
\end{solution}
